\documentclass[12pt]{article}
\usepackage[utf8]{inputenc}
\usepackage[T1]{fontenc}
\usepackage[margin=1in]{geometry}
\usepackage{setspace}
\usepackage{xcolor}
\usepackage[colorlinks=true, urlcolor=blue!60!black, linkcolor=blue!60!black]{hyperref}
\setlength{\parindent}{2em}
\setlength{\parskip}{1.25em}
\renewcommand{\baselinestretch}{1.0}
\begin{document}
\begin{center}
{\LARGE \textit{Wickard v. Filburn}, 317 U.S. 111 (1942)\par}
\vspace{0.8em}
{\large Case Brief --- Four-Page Edition\par}
\end{center}
\vspace{1em}

\section{Memory Jogger}

Under the Commerce Clause, Congress may regulate purely local, non-marketed economic activity, even a farmer's home-grown wheat consumed on his own farm, when that activity, aggregated with the conduct of many others similarly situated, exerts a substantial economic effect on interstate commerce (\href{https://supreme.justia.com/cases/federal/us/317/111/}{317 U.S. 111}).

\section{Detailed Case Facts}

Roscoe Filburn ran a small farm in Montgomery County, Ohio, keeping dairy cattle and poultry. Each fall he sowed winter wheat, selling part of the crop, feeding part to his animals, grinding some into flour, and keeping the rest as seed.

To stabilize collapsing Depression-era prices, the Agricultural Adjustment Act of 1938 authorized national wheat acreage allotments, apportioned down to individual farms, and compulsory marketing quotas approved by two-thirds of farmers in a referendum. In July 1940, Filburn was allotted 11.1 acres but sowed 23. Shortly before the May 1941 referendum, Secretary of Agriculture Claude Wickard gave a radio address urging approval without mentioning a pending amendment that raised the penalty on excess wheat from 15 to 49 cents a bushel and extended it to wheat consumed on the farm. Farmers approved the quota, 81 to 19 percent. Filburn's 239 excess bushels drew a \$117.11 penalty, which he refused to pay. He sued, arguing that Congress could not regulate wheat grown for home use and that the amendment was impermissibly retroactive.

\section{Procedural History}

A divided three-judge district court enjoined enforcement, holding that the radio speech tainted the referendum and that the penalty increase was retroactive. On direct appeal, the Supreme Court heard argument in May 1942, ordered reargument on the Commerce Clause question, and reversed on November 9, 1942.

\section{Judicial Votes}

Reversed, \textbf{9--0}.
\begin{itemize}
\item \textbf{Opinion of the Court:} Justice Jackson, joined by Chief Justice Stone and Justices Roberts, Black, Reed, Frankfurter, Douglas, Murphy, and Byrnes.
\item \textbf{Concurring or dissenting:} None.
\end{itemize}

\section{Holding}

The commerce power reaches local activity that is not itself commerce, including production for home consumption, if it exerts a substantial economic effect on interstate commerce, regardless of whether the effect would once have been called direct or indirect. Under the aggregation principle, a trivial individual contribution does not escape regulation when the combined effect of many similarly situated actors is substantial. The penalty was not unconstitutionally retroactive, and the radio speech did not invalidate the referendum. Jackson's suggestion that restraints on the commerce power come mainly from politics went beyond the case; \textit{Lopez}, \textit{Morrison}, and \textit{NFIB} later reaffirmed judicial limits.

\section{Analysis of Opinions}

\textbf{Opinion of the Court (Jackson).} Jackson dismantled the pre-1937 categories that treated production and manufacturing as local and their effects as indirect as a matter of law, as in \textit{E.\ C.\ Knight} and \textit{Carter Coal}. After rejecting the district court's referendum ruling as ``manifest error,'' he traced doctrine from \textit{Gibbons v.\ Ogden} through \textit{Swift} and the \textit{Shreveport Rate Cases}, replacing labels with a practical inquiry into economic effects. Whether Filburn's conduct was called production, consumption, or marketing did not matter.

The stipulated record showed an integrated national and world wheat market, with domestic prices supported at about \$1.16 a bushel against a 40-cent world price. Home-consumed wheat was the most variable factor in wheat disappearance. It affected prices in two ways: excess wheat overhangs the market and can flow into commerce if prices rise, and even unsold wheat satisfies the grower's own needs, forestalling purchases on the open market. Filburn's 239 bushels were trivial alone, but the aggregate of thousands of farmers was not, and conflicts between small and commercial growers were for Congress.

On due process, the scheme was not arbitrary, particularly since cooperating farmers received price supports; it is hardly a denial of due process to regulate what the government subsidizes. Nor was the amendment retroactive, because the penalty attached only on threshing after enactment, and Filburn could have stored the excess or fed it unthreshed.

\section{Examples: Future Applications}

\begin{enumerate}
\item \textbf{Same-Side (Valid): Home-Grown Marijuana.} A medical-cannabis patient grows six plants for personal use. Following \textit{Wickard}, \textit{Gonzales v.\ Raich} upheld federal regulation because home-grown fungible commodities affect interstate supply and demand in the aggregate.
\item \textbf{Same-Side (Valid): Backyard Smelting.} A federal scheme caps backyard coal smelting, reaching a workshop that smelts scrap only for its own tools. Production of a fungible, traded commodity can be aggregated.
\item \textbf{Opposite-Side (Invalid): Guns Near Schools.} \textit{Lopez} and \textit{Morrison} held that non-economic possession or violence cannot be aggregated into a general police power.
\item \textbf{Opposite-Side (Invalid): A Produce Mandate.} Requiring non-growers to buy \$200 of produce each year compels commerce, which \textit{NFIB v.\ Sebelius} forbids; Filburn was already farming.
\item \textbf{Fence-Sitter: Home Distilling.} A ban on distilling spirits for personal use protects tax revenue and a national market, but challengers call hobby distilling non-commercial, and some courts have recently questioned the ban.
\end{enumerate}

\section{Critique}

\textbf{Originalist and textualist views.} Scholars Randy Barnett and Richard Epstein argue that \textit{Wickard} departed from the original meaning of regulating commerce among the states, which referred to trade across state lines rather than local agriculture or home consumption. Pairing substantial effects with aggregation, critics say, collapses enumerated powers and erodes the Tenth Amendment, since nearly any activity affects the economy when multiplied by millions. Others note that the decision rests implicitly on a broad reading of the Necessary and Proper Clause.

\textbf{Progressive and pragmatic views.} Defenders see \textit{Wickard} as a necessary adaptation to an integrated economy and to collective-action problems, like price collapse and overproduction, that states cannot solve alone. Its rejection of the direct-indirect formalism underlies federal workplace laws and the Civil Rights Act of 1964, while \textit{Lopez}, \textit{Morrison}, and \textit{NFIB} show that courts can keep the aggregation rule without granting unlimited power.

\section{Key Quotations}

\begin{quote}
``But even if appellee's activity be local, and though it may not be regarded as commerce, it may still, whatever its nature, be reached by Congress if it exerts a substantial economic effect on interstate commerce, and this irrespective of whether such effect is what might at some earlier time have been defined as `direct' or `indirect.' ''
\end{quote}
\noindent --- Justice Jackson, Opinion of the Court, 317 U.S. at 125.

\begin{quote}
``That appellee's own contribution to the demand for wheat may be trivial by itself is not enough to remove him from the scope of federal regulation where, as here, his contribution, taken together with that of many others similarly situated, is far from trivial.''
\end{quote}
\noindent --- Justice Jackson, Opinion of the Court, 317 U.S. at 127--128.

\begin{quote}
``Home-grown wheat in this sense competes with wheat in commerce. The stimulation of commerce is a use of the regulatory function quite as definitely as prohibitions or restrictions thereon.''
\end{quote}
\noindent --- Justice Jackson, Opinion of the Court, 317 U.S. at 128.

\end{document}
